% \section{Flasher un ZYNQ avec l'infrastructure FPGA}

\subsection{Generating the encrypted bitstream}

The default behavious of Vivado is to generate a bistream file with the {\tt .bit} 
extension. The FPGA Manager expects another format with the appropriate header. Converting
from {\tt .bit} to {\tt .bit.bin} is achieved with the {\tt bootgen} tool provided with
the Vivado SDK.

This tool expects a configuration file (called here {\tt filename.bif}) with the {\tt .bif} 
extension as simple as: 
\begin{lstlisting}
all:
{
  bitstream_name.bit
}
\end{lstlisting}

which is provided as argument to {\tt bootgen}:
\begin{lstlisting}[language=bash]
$VIVADO_SDK/bin/bootgen -image filename.bif -arch zynq -process_bitstream bin
\end{lstlisting}

Following the execution of this command, a file named {\tt bitstream\_name.bit.bin} has been
created in the current directory.

\subsection{Configuring the FPGA through {\tt fpga\_manager}}

The {\tt .bit.bin} file must be copied/moved to the {\tt /lib/firmware} directory.

In order to tell the driver that the {\tt PL} must be configured and which bitstream
is to be used, the following command is used:
\begin{lstlisting}[language=bash]
echo "bitstream_name.bit.bin" > /sys/class/fpga_manager/fpga0/firmware
\end{lstlisting}
leading to the display of the line
\begin{lstlisting}[language=bash]
fpga_manager fpga0: writing nom_du_bitstream.bit.bin to Xilinx Zynq FPGA Manager
\end{lstlisting}
in the console or the kernel log, and the LED connected to {\tt Prog done} must
switch on (blue LED on the Red Pitaya).

\subsection{Using a devicetree overlay for configuring the FPGA}

As was the case with the previous solution, the bitstream must be located in
{\tt /lib/firmware}.

The benefit of this method over the previous one is that when the design includes IPs
communicating with the processor, the {\em devicetree} tells the kernel which drivers
to load. This driver is provided as a kernel module to be inserted when the bitstream
is used to configure the FPGA. The overlay overloads the {\tt fpga\_full} node definition 
by providing simultaneously the bistream name as well as the identifier of the dedicated
driver to be loaded for this application.

Without getting in all the details of the overlay format, the following example
will allow to:
\begin{itemize}
\item modify the {\tt fpga\_full} node definition, {\tt target} attribute, telling
{\tt fpga\_manager} the name of the binary file to be used for configuring the
PL through the {\tt firmware-name} attribute;
\item add a node defining the drivers needed to control the PL from the PS leading
to the module being loaded. In the current example the associated driver is
called {\tt gpio\_ctl} ({\tt compatible} field), and the base address for sharing
information between PS and PL (0x43C00000) is provided, as well as the size (0x1f) 
through the {\tt reg} attribute.
\end{itemize}\vspace{0.3cm}
 
\begin{lstlisting}
/dts-v1/;
/plugin/;
/ {
    compatible = "xlnx,zynq-7000";
    fragment@0 {
        target = <&fpga_full>;
        #address-cells = <1>;
        #size-cells = <1>;
        __overlay__ {
            #address-cells = <1>;
            #size-cells = <1>;

            firmware-name = "top_redpitaya_axi_gpio_ctl.bin";

            gpio1: gpio@43C00000 {
                compatible = "gpio_ctl";
                reg = <0x43C00000 0x0001f>;
                gpio-controller;
                #gpio-cells = <1>;
                ngpio= <8>;
            };
        };
    };
};

\end{lstlisting}

This file must be compiled using the following command:
\begin{lstlisting}
/somewhere/buildroot/output/host/usr/bin/dtc -@ -I dts -O dtb -o ${FILENAME}.dtbo ${FILENAME}.dts
\end{lstlisting}
\noindent with:
\begin{itemize}
\item {\tt -@} to generate symbols to be dynamically loaded when 
inserting the module
\item {\tt -I dts} to define the kind of input file (DeviceTree Source);
\item {\tt -O dtb} to define the kind of output file (DevideTree Binary);
\item {\tt -o} the output filename.
\end{itemize}

Loading the file into the memory requires two steps:
\begin{enumerate}
\item creating a directory that will hold the overlay:
\begin{lstlisting}[language=bash]
mkdir /sys/kernel/config/device-tree/overlays/toto
\end{lstlisting}
will automaticall create a whole set of files in this directory:
\begin{lstlisting}[language=bash]
redpitaya> ls -l /sys/kernel/config/device-tree/overlays/toto/
total 0
-rw-r--r--    1 root     root             0 Jan  1 00:04 dtbo
-rw-r--r--    1 root     root          4096 Jan  1 00:04 path
-r--r--r--    1 root     root          4096 Jan  1 00:04 status
\end{lstlisting}
\item loading the overlay in the {\em devicetree}:
\begin{lstlisting}[language=bash]
cat gpio_red.dtbo > /sys/kernel/config/device-tree/overlays/toto/dtbo
\end{lstlisting}
will configure the FPGA by transfering the bitstram, load the driver by inserting the associated
module identified through the ``compatible'' argument which must match the corresponding
field in the kernel module.
\end{enumerate}

The modifications induced by the overlay can be canceled to return to the initial
devicetree state by deleting the directory holding the overlay: 
\begin{lstlisting}[language=bash]
rmdir /sys/kernel/config/device-tree/overlays/toto
\end{lstlisting}
